\chapter{Capstone 课程日历与里程碑}

\section{案例目标}

第 43 章把 Capstone 材料放进试教场景里检查。本章继续把它推进到课程层面：\textbf{怎样把 proposal、pilot、project board、rubric、报告、notebook freeze 和最终展示安排进一个 12 到 16 周的课程日历。}

本章的目标有四点：

\begin{enumerate}
    \item 理解课程日历不是日期列表，而是一个降低期末风险的项目管理工具；
    \item 学会把学生交付物、教师检查点、AI 使用规范和风险状态放到同一张周历里；
    \item 学会区分普通教学周、checkpoint week、instructor intervention week 和 final delivery week；
    \item 学会用课程日历提前发现“第 15 周才暴露”的范围、数据、复现和签字问题。
\end{enumerate}

\section{好的课程日历会把风险提前暴露}

很多 Capstone 课程最危险的地方，不是学生第 15 周还没写完报告，而是教师第 15 周才第一次发现：

\begin{itemize}
    \item 数据许可证一直没有确认；
    \item baseline 从未真正跑通；
    \item notebook 只能在一台电脑上运行；
    \item 学生一直把模型分数当成唯一结果；
    \item AI 使用声明到最后才补；
    \item 助教和教师对同一项目的评分尺度不一致。
\end{itemize}

所以，本章的核心立场是：\textbf{课程日历的任务不是把工作拖到截止日期，而是把关键风险尽早变成可见 checkpoint。}

\section{一个极小的 course-calendar 数据集}

配套 notebook 使用 \nolinkurl{capstone_course_calendar_demo.csv}。这是一组完全本地、教学用途的\textbf{Capstone 课程周历卡片}。每一行代表一周，包含：

\begin{itemize}
    \item \texttt{week}；
    \item \texttt{phase}；
    \item \texttt{milestone}；
    \item \texttt{student\_deliverable}；
    \item \texttt{instructor\_check}；
    \item \texttt{risk\_status}；
    \item \texttt{ai\_policy\_focus}；
    \item \texttt{workload\_score}；
    \item \texttt{reference\_route}。
\end{itemize}

这里的 \texttt{reference\_route} 分成四类：

\begin{itemize}
    \item \texttt{regular\_instruction}：普通教学和项目推进周；
    \item \texttt{checkpoint\_week}：学生有明确交付物或教师有检查项；
    \item \texttt{instructor\_intervention}：范围、数据边界或项目风险需要教师主动介入；
    \item \texttt{final\_delivery\_week}：最终报告、展示、归档和签字周。
\end{itemize}

这个数据集采用 16 周版本，但它可以压缩到 12 周。压缩时不要简单删掉 checkpoint，而应该合并普通教学周，保留 proposal、pilot、midpoint review、notebook freeze 和 final delivery 这些风险控制点。

\section{先看一个危险 baseline：只看 workload}

课程排期时，一个常见 baseline 是只盯着 workload score：哪一周任务重，哪一周就标成高风险；任务不重，就当作普通周。

\begin{examplebox}[只按 workload 粗略判断课程周状态]
\begin{lstlisting}[style=pythonstyle]
if row["workload_score"] > 0.70:
    intervention_week
elif row["week"] >= 15:
    final_delivery_week
else:
    regular_instruction
\end{lstlisting}
\end{examplebox}

这个 baseline 会漏掉很多真正重要的 checkpoint。例如第 2 周的数据边界、第 4 周的 one-week pilot、第 12 周的 notebook freeze，工作量不一定最高，但它们决定了项目后期是否可控。

在本章教学数据上，只按 workload 和周数粗略判断，route accuracy 只有 0.25。表~\ref{tab:ch44_workload_vs_calendar} 展示了结构化 course-calendar workflow 如何先看风险状态，再看最终交付周，最后检查学生交付物和教师检查项：

\begin{table}[htbp]
\centering
\small
\setlength{\tabcolsep}{4pt}
\begin{tabular}{@{}p{0.28\textwidth}>{\centering\arraybackslash}p{0.18\textwidth}p{0.44\textwidth}@{}}
\toprule
方法 & 周状态准确率 & 教学解读 \\
\midrule
只看 workload 和周数 & 0.25 & 会漏掉很多低负担但高价值的 checkpoint \\
结构化 course-calendar workflow & 1.00 & 先识别风险和最终交付，再安排学生与教师检查点 \\
\bottomrule
\end{tabular}
\caption{本章 notebook 中两个最小课程排期流程的对照。课程日历的重点不是平均分配工作量，而是提前暴露风险。}
\label{tab:ch44_workload_vs_calendar}
\end{table}

\section{一个可执行的 course-calendar workflow}

本章 notebook 的 routing 逻辑可以写得很小：

\begin{examplebox}[一个最小课程周历 routing 逻辑]
\begin{lstlisting}[style=pythonstyle]
if risk_status == "high":
    instructor_intervention
elif phase in {"final_delivery", "showcase"}:
    final_delivery_week
elif student_deliverable != "none":
    checkpoint_week
elif instructor_check != "none":
    checkpoint_week
else:
    regular_instruction
\end{lstlisting}
\end{examplebox}

这段逻辑表达了一个课程设计原则：\textbf{风险优先于日期，交付优先于讲授。} 如果某一周风险已经高，就不能只把它当作普通教学周；如果某一周有学生交付物或教师检查项，就应该明确写进 calendar，而不是口头提醒。

\section{一份 16 周 Capstone 日历的骨架}

一个可直接改写的 16 周骨架可以是：

\begin{enumerate}
    \item Kickoff：课程目标、环境检查、AI 使用边界；
    \item Proposal card：科学问题、数据来源和公开边界；
    \item Data inventory：字段、许可证、清洗规则和 notebook 路径；
    \item One-week pilot：最小 baseline 和第一条失败样本；
    \item Baseline result：指标 sanity check 和简单报告；
    \item Model comparison：尝试改进模型，但不扩大范围；
    \item Claim ledger：文献、claim、caveat 和引用核查；
    \item Midpoint project board：教师介入处理范围漂移；
    \item Validation report：验证集、测试集、数据泄漏检查；
    \item Error analysis：失败样本、局限和不确定度；
    \item Report outline：报告结构和 AI-use statement 草稿；
    \item Notebook freeze：从头复跑、环境记录和随机种子；
    \item Scope lock：停止扩大题目，处理最终风险；
    \item Presentation rehearsal：展示计时和核心图选择；
    \item Final package：报告、notebook、图表、声明和签字；
    \item Showcase and archive：展示、归档和公开边界检查。
\end{enumerate}

如果课程只有 12 周，可以合并第 5 到第 7 周、压缩第 9 到第 11 周，但不要删掉第 4 周 pilot、第 8 周 midpoint review、第 12 周 notebook freeze 和最终交付周。

\section{每个 checkpoint 都应该有一个出口}

课程日历最容易写成“本周讲什么”。但对 Capstone 来说，更有用的写法是“本周结束时项目应该处于什么状态”。例如：

\begin{itemize}
    \item 第 2 周结束：每个项目都有 proposal card 和数据边界；
    \item 第 4 周结束：每个项目至少跑通一个 baseline；
    \item 第 8 周结束：所有项目都经过 project board review；
    \item 第 12 周结束：notebook 能从头复跑；
    \item 第 15 周结束：报告、展示、AI-use statement 和 signoff 都进入最终检查。
\end{itemize}

这让课程日历变成一组状态门，而不是一串日期。

\section{配套 notebook 做了什么}

本章 notebook 采用的是一个 teaching-first 的最小设计：

\begin{itemize}
    \item 先读取 course-calendar table，并统计 phase、risk 和 reference route；
    \item 实现一个只按 workload 和周数判断的 naive baseline；
    \item 再实现一个带 risk、final delivery、student deliverable 和 instructor check 的 calendar workflow；
    \item 输出 regular、checkpoint、intervention 和 final delivery 四个队列；
    \item 为 intervention 和 final delivery 周生成课程团队 checklist；
    \item 最后生成一个 12 周压缩版建议和 calendar assistant prompt。
\end{itemize}

这份 notebook 不追求复杂排课软件，而是给课程团队一个可检查的、能复用的 Capstone 周历骨架。

\section{AI 助手如何帮助你}

在这个案例里，AI 助手最适合帮助你：

\begin{itemize}
    \item 把 16 周日历压缩成 12 周版本，并说明哪些 checkpoint 不能删；
    \item 检查每一周是否有清楚的学生交付物或教师检查项；
    \item 把试教反馈映射到下一轮课程日历修订；
    \item 帮你生成学生版日历、助教版日历和教师内部风险表；
    \item 检查 AI 使用规范是否在关键周重复出现，而不是只在第一周讲一次。
\end{itemize}

但课程节奏、评分政策、工作量和校历限制仍然需要课程团队自己判断。AI 助手可以给出草案，不能替代真实学生负担和学校制度。

\section{练习}

\begin{exercisebox}[基础练习]
把第 12 周的 \texttt{instructor\_check} 改成 \texttt{none}。重新运行 notebook，观察 notebook freeze 为什么不应该被当成普通教学周。
\end{exercisebox}

\begin{exercisebox}[进阶练习]
把课程压缩到 12 周。请说明你会合并哪些普通教学周，以及为什么仍然保留 proposal、pilot、midpoint review、notebook freeze 和 final delivery。
\end{exercisebox}

\begin{exercisebox}[开放练习]
为你自己的课程项目写一份周历草案。每周至少包含：phase、学生交付物、教师检查项、AI 使用规范焦点和风险状态。
\end{exercisebox}

\section{本章小结}

Capstone 课程日历不是行政排期，而是项目风险控制工具。本章把 phase、milestone、学生交付物、教师检查项、AI 使用规范、workload 和风险状态放进同一张 course-calendar card。

当课程团队能在第 2 周看到数据边界、第 4 周看到 baseline、第 8 周看到 project board、第 12 周看到 notebook freeze，期末交付就不再是一场惊险冲刺，而是一条可复盘的科研训练路径。
